\documentclass[11pt]{amsart}
\usepackage[T1]{fontenc}
\usepackage{lmodern}
\usepackage{amsmath,amssymb,amsthm}
\usepackage{xcolor}
\usepackage[colorlinks=true,linkcolor=blue!55!black,citecolor=blue!55!black,urlcolor=blue!55!black]{hyperref}

\newtheorem{theorem}{Theorem}
\newtheorem{lemma}[theorem]{Lemma}
\newtheorem{remark}[theorem]{Remark}
\title{Strict Domination Growth from $J(8,4)$ to $J(9,4)$}
\author{Xiaojian Zeng}
\address{PBC School of Finance, Tsinghua University, 100190}
\email{xiaojianzeng@foxmail.com}
\author{Jianming Wang}
\address{Shanghai Siqiang Intelligent Control Robot Engineering Co., Ltd}
\email{847017136@qq.com}
\date{September 2, 2026}
\hypersetup{pdftitle={Strict Domination Growth from J(8,4) to J(9,4)},pdfauthor={Xiaojian Zeng; Jianming Wang}}

\begin{document}
\begin{abstract}
We prove $\gamma(J(8,4))<\gamma(J(9,4))$, the first parameter pair in the
$r=4$ case of the strict-monotonicity conjecture for domination numbers of
Johnson graphs.  An explicit six-vertex family dominates $J(8,4)$.  On the
other hand, a six-vertex dominating family in $J(9,4)$ would partition the
graph into radius-one balls and thereby induce a $2$-$(9,4,1)$ design.  The
latter is impossible by the point-replication equation $3r_x=8$.  Thus the
proof needs neither exact domination numbers nor a solver optimality claim.
\end{abstract}
\maketitle

\section{Introduction}
For positive integers $n\ge r$, the Johnson graph $J(n,r)$ has vertex set
$\binom{[n]}r$, with two distinct vertices adjacent when their intersection
has size $r-1$.  A set $D$ of vertices dominates a graph when every vertex is
in $D$ or adjacent to a member of $D$; the minimum possible size is denoted by
$\gamma$.

Cornet, Dravec, and Torres conjectured that
\[
 \gamma(J(n,r))<\gamma(J(n+1,r))\qquad(r\ge3,\ n\ge2r)
\]
\cite[Conjecture~14]{CornetDravecTorres2025}.  Together with the even-$n$
result of Cornet, Dravec, and Torres, Lee and Oh completed the case $r=3$
and recorded in June 2026 that the cases $r\ge4$ remained open
\cite[Concluding remarks]{LeeOh2026}.  We settle the boundary parameter pair
$(r,n)=(4,8)$.

\begin{theorem}\label{thm:main}
One has
\[
 \gamma(J(8,4))\le6<7\le\gamma(J(9,4)).
\]
In particular, $\gamma(J(8,4))<\gamma(J(9,4))$.
\end{theorem}

The theorem intentionally does not assert either domination number exactly.
The separated bounds are enough, and their proof is elementary apart from a
small explicit witness.

\section{A radius-one obstruction}
For $A\in\binom{[n]}4$, write
\[
 B(A)=\{X\in\tbinom{[n]}4:|A\cap X|\ge3\}
\]
for its closed neighbourhood.  Replacing no element or exactly one element of
$A$ gives
\begin{equation}\label{eq:ballsize}
 |B(A)|=1+4(n-4).
\end{equation}

\begin{lemma}\label{lem:intersection}
If $A,C\in\binom{[n]}4$ and $|A\cap C|\ge2$, then
$B(A)\cap B(C)\ne\varnothing$.
\end{lemma}
\begin{proof}
If $|A\cap C|\ge3$, one of the centres itself belongs to both closed
neighbourhoods.  If $A\cap C=Q$ has size two, choose
$a\in A\setminus C$ and $c\in C\setminus A$.  Then the four-set
$Q\cup\{a,c\}$ meets each centre in three points.
\end{proof}

\begin{lemma}\label{lem:lower}
$\gamma(J(9,4))\ge7$.
\end{lemma}
\begin{proof}
There are $\binom94=126$ vertices, while \eqref{eq:ballsize} gives
$|B(A)|=21$.  Fewer than six balls cannot cover all vertices.  If six balls
did cover them, equality in
\[
 \left|\bigcup_{A\in D}B(A)\right|\le\sum_{A\in D}|B(A)|=6\cdot21=126
\]
would force the balls to be pairwise disjoint.  Lemma~\ref{lem:intersection}
then implies that the six centre four-sets meet pairwise in at most one point.

Consequently, the six centres contain $6\binom42=36$ distinct unordered
pairs.  Since $\binom92=36$, every pair of points of $[9]$ occurs in exactly
one centre.  Fix a point $x$ and let $r_x$ be the number of centres containing
$x$.  Each such centre accounts for three of the eight pairs containing $x$,
and every such pair is accounted for exactly once.  Hence $3r_x=8$, an
impossibility.  Thus six balls cannot dominate, proving the claim.
\end{proof}

\section{The explicit upper bound}
\begin{lemma}\label{lem:upper}
$\gamma(J(8,4))\le6$.
\end{lemma}
\begin{proof}
Consider the following family:
\begin{align*}
\mathcal D=\{&1256,\ 1348,\ 1678,\\
             &2347,\ 2568,\ 3457\},
\end{align*}
where, for example, $1256$ abbreviates $\{1,2,5,6\}$.  Direct enumeration of
the $\binom84=70$ four-sets shows that each meets a member of $\mathcal D$ in
at least three points.  For completeness, the accompanying verifier performs
this check twice: once by the intersection predicate and once by explicitly
generating all one-element replacements.  Therefore $\mathcal D$ dominates
$J(8,4)$.
\end{proof}

\begin{proof}[Proof of Theorem~\ref{thm:main}]
Combine Lemmas~\ref{lem:lower} and~\ref{lem:upper}.
\end{proof}

\section{Reproducibility and limitations}
The witness is stored as a literal JSON list, and a short Python verifier
checks its syntax, both evaluations of its coverage, all closed-ball sizes,
and the local intersection lemma for every pair of centres in $J(9,4)$.  The
command is \texttt{python3 code/verify\_certificate.py}.  All arithmetic is
integer and deterministic.  CP-SAT was used only to discover the six-set
witness; no solver status, floating-point value, or uncheckable search output
is used in the proof.

The result proves only the strict inequality at $(n,r)=(8,4)$.  It neither
determines the two domination numbers individually nor addresses later
$r=4$ instances.  Literature and priority searches were conducted through
September 2, 2026.  The mathematical certificate and every logical
implication needed for the theorem are included in the source package for
independent checking.

\bibliographystyle{amsplain}
\bibliography{references}
\end{document}
